\chapter{Rationality and the Public Understanding of Science}
\label{sec:p18-rationality-14-3}

A person can understand a scientist's conclusion without being able to
evaluate the instruments, models and evidence that warrant it. The same
person may nevertheless have reasons to prefer that scientist's account
to a competing source. Bromme and Gierth's 2021 chapter examines this
position: public understanding of science is important for everyday
rationality, yet individual understanding of specialized science remains
bounded. Their central argument is that assessing \emph{whom to believe}
can be a rational response to those bounds, not a renunciation of reasons
\cite{BrommeGierth2021PublicScience}.

The treatment below reconstructs the complete chapter's argument, examples
and qualifications. A bounded original empirical record uses Hendriks,
Kienhues and Bromme's 2015 Muenster Epistemic Trustworthiness Inventory
study \cite{HendriksKienhuesBromme2015METI}. An inventory of perceived
trustworthiness is not a test of the actual reliability of science, and
neither perceived reliability nor scientific knowledge settles every
public decision.

\section{Science, reasoning standards and background knowledge}
\label{sec:p18-r143-nexus}

\begin{hbdefinition}{def:p18-r143-scope}{Science and two senses of public understanding}
The source uses \emph{science} in the broad sense of \emph{Wissenschaft},
including social sciences and humanities, not only experimental natural
science. Its \emph{analytic} concept of public understanding of science
embraces citizens' encounters with scientific knowledge, claims, institutions
and controversies. A \emph{strategic} use of the expression concerns efforts
to educate, communicate or promote appreciation and acceptance of science.
Describing an encounter is not the same as endorsing a campaign's goals or
establishing its effectiveness.
\end{hbdefinition}

The source's nexus is reciprocal. Scientific production and justification
must be answerable to standards of rational inquiry; conversely, many
everyday inferences depend on background knowledge produced through
scientific inquiry. A coherent calculation on false or inapplicable
background assumptions need not describe the world. Public access to
methods, models and their warranted conclusions therefore matters, not
just adherence to an inference pattern.

The chapter illustrates accepted reasoning standards and background
knowledge using Bayesian inference and Newton's first law. These have
different kinds of warrant: a probability identity follows under formal
assumptions, whereas applying a physical law additionally requires an
appropriate model and evidence. Newton's constant-velocity example
presupposes an inertial frame and zero net force; it is not an unconditional
statement that every object keeps the same velocity. Community acceptance
alone supplies neither a mathematical proof nor those physical conditions.
The authors' point is the role of accepted knowledge in evaluating an
inference, not the identity of mathematical and empirical justification
\cite[pp.~767--768]{BrommeGierth2021PublicScience}.

Science and rationality are not equivalent. Rational reasoning need not
refer to science in every instance, and being locally rational is not
sufficient to make a practice scientific. Note~3 explicitly rejects the
claim that only scientific propositions can be rationally warranted and
recognizes limits of scientific explanation. Thus the argument is not a
scientistic monopoly on reasons. Note~4 contrasts truth-directed beliefs
about the world's structure with arbitrarily adopted beliefs about aliens:
local consistency does not remove the demand for warranted background
knowledge \cite[pp.~767, 773, notes~3--4]{BrommeGierth2021PublicScience}.

The demarcation discussion illustrates why local coherence is insufficient.
The chapter names creationism, climate-change denialism, homeopathy,
flat-earth and chemtrail theories as examples of pseudoscientific frameworks.
Such frameworks can imitate scientific discourse and offer purported
evidence. They are not identical to every conspiracy belief or superstition;
the source recognizes fuzzy boundaries. Nor does it claim a single precise
set of necessary and sufficient demarcation conditions. Its proposed
assessment instead attends to neglected scientific standards and conflicting
evidence.

In the homeopathy example, an inference about increasing potency through
dilution can fit an assumed world in which substances preserve health-related
information independently of molecular concentration. Consistency with
that assumption does not establish its evidential warrant. Selectively
discarding contrary observations creates the broader epistemic problem.
This is an illustration of the difference between conditional reasoning
and justified premises, not medical advice or a clinical experiment
performed by this chapter \cite[p.~767]{BrommeGierth2021PublicScience}.

\section{Scientific literacy and bounded public understanding}
\label{sec:p18-r143-boundedness}

The chapter distinguishes three research strands: scientific-literacy
education, especially for children and young adults; adults' processing
and acceptance of scientific claims in media; and population surveys of
stances toward science. Their tasks, populations and measures are not
interchangeable. A science-test score, an attitude rating and a decision
to rely on an expert supply different evidence.

\begin{hbdefinition}{def:p18-r143-literacy}{Literacy and encounters beyond the curriculum}
In the source's OECD-based account, scientific literacy involves explaining
phenomena, evaluating or designing inquiry, and interpreting data,
evidence and arguments. School-oriented literacy assessments use selected
knowledge and reasoning accessible without specialist training.
Public understanding also concerns unfamiliar, uncertain or contested
claims encountered for personal or collective purposes. It is not confined
to knowledge suitable for inclusion in an introductory curriculum.

\emph{Bounded understanding} means limited ability to understand or assess
the specialized content and its justification. It need not mean an empty
mind, an immutable inability, a general intellectual defect, or an
automatic cause of rejecting science.
\end{hbdefinition}

Parents assessing educational practices, patients comparing a clinician's
account with internet information, and citizens considering carbon storage
as a climate response are the source's examples. What attracts public
attention depends on practical questions and media coverage, not only
on what is easy to understand or already settled. Scientific-literacy
differences observed in PISA cannot by themselves characterize every
adult encountering these topics \cite[p.~768]{BrommeGierth2021PublicScience}.

The source develops several reasons for boundedness. Explanations can
extend through many levels; evidence production involves technological,
social and epistemic complexity; disciplines divide cognitive work.
Some entities, such as genes, viruses and bacteria, cannot be observed
with unaided senses. Mathematical models and scientific categories may
depart substantially from ordinary experience. An accessible analogy
can assist comprehension without conveying every condition of the model.

Everyday understanding also contains stable, locally coherent models,
not merely missing facts. The chapter contrasts object/substance-like
explanations with relational scientific models, and essence-like genetic
explanations with more complex accounts of development and behavior.
Mass, energy and heat are the authors' physics illustrations of that
contrast. Their distinct scientific definitions are not reduced here
to one common relational axiom.
Those contrasts identify possible conceptual obstacles; they are not
a rule that every everyday explanation is wrong or every scientific
concept belongs to a single ontology. The reviewed misconception evidence
has no single common participant denominator or effect estimate in this
chapter \cite[pp.~768--769]{BrommeGierth2021PublicScience}.

The boundary can move. Nuclear physics and molecular genetics illustrate
specialist concepts entering public education; instruction and well-designed
communication can improve understanding. This does not imply that all
specialist evidence becomes independently assessable by every citizen.
Disagreement creates an additional problem: understanding the language
of two incompatible claims does not establish which is warranted.
The source reports that 70 percent of respondents in Castell and colleagues'
2014 British survey agreed that conflicting information makes it hard
to know what to believe. The chapter does not provide that item's sample
size, response scale, weighting or interval; the number remains a
review-reported survey result, not a universal public-confusion rate
\cite[p.~769 and note~5]{BrommeGierth2021PublicScience}.

\section{From claim plausibility to source assessment}
\label{sec:p18-r143-trust}

\begin{hbdefinition}{def:p18-r143-epistemic-trust}{Dependence, vulnerability and epistemic trust}
Trust involves willingness to depend on another agent while accepting
risks associated with that dependence. A trustor relies on a trustee,
which can be a person or an institution. In \emph{epistemic} trust the
relevant good is warranted knowledge, and the vulnerability concerns
error, misleading information or unjustified belief.

\emph{Perceived trustworthiness} is an assessment of the source.
\emph{Actual reliability} concerns how its methods or reports track the
relevant truth under the conditions at issue. \emph{Deference} is reliance
on the source, potentially expressed in belief or action. These notions
are related but do not define each other into equivalence.
\end{hbdefinition}

Judging content asks whether a claim fits available evidence and knowledge.
Judging a source asks whether its relevant competence, conduct and
institutional setting justify reliance. Their targets differ even when
each informs the other. A person's disagreement with a message can lower
their assessment of its source; their initial confidence in a source
can increase the message's plausibility. The interaction can supply
useful evidence or reinforce a motivated circle. It must be examined
rather than treated as independent confirmation.

Unqualified authority reasoning does not establish deductive truth:
a genuinely knowledgeable expert can make a mistake. Yet a defeasible
argument from relevant expert opinion can be rational. Rejecting every
such argument because it is not deductively conclusive would also reject
ordinary reliance on distributed expertise. The source locates its
argument in that distinction: asking whom to believe changes the subject
of assessment, rather than abolishing assessment
\cite[pp.~769--770]{BrommeGierth2021PublicScience}.

The division of cognitive labor supplies the practical premise. Scientific
specialists themselves cannot independently reproduce every observation,
instrument calibration and calculation on which their work depends.
Citizens face a still broader set of specialized claims. A demand for
complete first-hand verification would therefore remove useful knowledge
without supplying an alternative way to obtain it. Reliance remains
answerable to qualifications and contrary evidence; it does not make
scientists infallible.

\begin{hbdefinition}{def:p18-r143-trust-dimensions}{Three dimensions and the pertinence condition}
\emph{Expertise} concerns knowledge and skill. \emph{Pertinence} asks
whether that expertise is relevant to this particular claim; professional
status in one domain does not confer competence in every other domain.
\emph{Integrity} concerns adherence to professional/scientific standards,
including honest, unbiased inquiry and reporting. \emph{Benevolence}
concerns regard for the trustor's or society's interests.

An expert can pursue personal interests while also acting benevolently.
A disclosed interest is relevant information, not automatic proof of
deception. Conversely, benevolent intentions cannot compensate by
themselves for a lack of relevant knowledge or defective methods.
\end{hbdefinition}

The chapter's source-processing review reports little spontaneous attention
to provenance and increased source consideration after discrepancies
within a text, between texts, or with prior beliefs. These are the
authors' broad literature summaries, not a universal law that every
person ignores sources until conflict. The chapter supplies no single
sample/design record or uncertainty estimate for that composite claim
\cite[p.~770]{BrommeGierth2021PublicScience}.

Its explanations of content-first judgment include overconfidence,
illusion of explanatory depth and processing fluency. In the reviewed
Scharrer, Stadtler and Bromme (2014) medical-text example, lexically easier
texts reduced reported need for expert advice despite difficult underlying
content; making epistemic complexity explicit mitigated the effect.
Comprehensibility is not competence to judge the evidence. That finding
remains review-level here: original participant counts, randomization,
materials, estimates and intervals are not reconstructed from the chapter
\cite[p.~771]{BrommeGierth2021PublicScience}.

\section{A primary record for measuring ascribed trustworthiness}
\label{sec:p18-r143-meti}

\begin{hbdefinition}{def:p18-r143-meti-record}{The three METI studies and what they measure}
Hendriks and colleagues developed a semantic-differential inventory of
ascribed expertise, integrity and benevolence toward an unfamiliar source.
Their fictional science-blog task described an experimental neuroenhancing
drug study and a researcher at a fictional institution. Names, institution
and medication were invented; these are not drug-efficacy data.

Study~1 obtained 237 completed surveys from 300 people opening its link.
Participants were volunteers previously registered for research invitations,
aged 19--47, mean 25.39, SD 4.38; 75.5 percent were women. Most had German
upper-secondary or university qualifications. Eighteen adjective-pair
items were initially presented on seven-point scales. Exploratory analysis
removed two, yielding sixteen items across three correlated factors.

Study~2 recruited students through lectures at a German university: 406
opened the survey and 345 completed it. Mean age was 21.77, SD 3.93;
69.3 percent were women. The same blog/procedure tested the measurement
structure in a new sample. Two more items were removed and a residual
covariance allowed, yielding fourteen items: six expertise, four integrity
and four benevolence. This is confirmation with post hoc modification,
not a wholly unchanged preregistered model.

Study~3 had 243 survey activations, 152 completions and 137 analysed people
after excluding the upper and lower five percent of completion times.
The retained window was 5.20--42.54 minutes. Mean age was 26.22, SD 5.24;
75.2 percent were women and 65 percent reported college/university
graduation. Each person rated all six source descriptions in randomized
order, crossing the targeted dimension with a high/low description.
The six responses per person are not six independently recruited samples.
\end{hbdefinition}

Completed surveys had no missing item values because of survey-software
settings, but incompletion still occurred before those analytic sets.
Voucher lotteries encouraged participation. The samples were predominantly
educated German volunteers/students, not representative national publics.
The article makes its data available through a repository and describes
supplementary materials; this record uses the main article, not a new
analysis of those participant data \cite[pp.~4--11]{HendriksKienhuesBromme2015METI}.

\begin{hbconvention}{def:p18-r143-statistical-notation}{Reported measurement statistics}
Here \(n\) is an analysed participant count, \(r\) a reported correlation,
\(\alpha\) the reported internal-consistency coefficient, \(F\) a
comparison statistic with its displayed degrees of freedom, and \(p\)
the source's reported significance level. The named fit indices describe
the measurement model's fit, not the probability that a source is truthful.
No quoted statistic is a fresh estimate from reanalysis.
\end{hbconvention}

Study~1's three-factor solution explained 61.66 percent of the item variance;
reported internal-consistency \(\alpha\) values were .91, .82 and .90 for
expertise, integrity and benevolence. Those values concern item relations,
not correctness of the researcher's claim. In Study~2, reported fit indices
were CFI .97, TLI .96, RMSEA .06 and SRMR .06. Multivariate-normality testing
led to use of a bootstrap. Integrity and benevolence correlated
\(r=.71\), while two integrity items had relatively low communalities.
The authors discuss overlapping aspects of character and values within
integrity. A correlated three-factor model is not three independent tests
of objective reliability \cite[pp.~5--9, 16]{HendriksKienhuesBromme2015METI}.

Study~3 supplied concrete cues. Low/high expertise compared a medical
student with an established neurology professor. Low integrity described
withholding findings contrary to earlier results; high integrity described
an open-science commitment. Low benevolence described prioritizing company
profit without adequate adverse-effect research; high benevolence described
aiming to help patients and society. All descriptions were presented as
relevant to neurology. These packages vary meaningful source descriptions,
not verified actual conduct of real researchers
\cite[pp.~9--11, table~5]{HendriksKienhuesBromme2015METI}.

The actual Table~6 reports the following mean ratings, with SDs in
parentheses. Its score orientation is higher for greater ascribed
trustworthiness; individual item endpoint directions and scoring must
not be replaced by a probability interpretation.
\begin{center}
\begin{tabular}{lrrr}
\toprule
Description & Expertise & Integrity & Benevolence\\
\midrule
Low expertise &3.80(.91)&4.62(.84)&4.45(.82)\\
High expertise &6.37(.88)&4.99(1.01)&5.01(1.07)\\
Low integrity &4.75(1.02)&2.40(.92)&2.51(1.00)\\
High integrity &5.90(.90)&5.93(.98)&5.79(1.01)\\
Low benevolence &4.74(1.08)&2.41(1.04)&1.88(.99)\\
High benevolence &5.91(.88)&5.67(.99)&5.87(1.02)\\
\bottomrule
\end{tabular}
\end{center}
The SDs are dispersion, not confidence intervals. The corresponding
descriptions affect more than their targeted subscale, so a change in
one perceived dimension cannot be interpreted as isolated from the others
\cite[p.~13]{HendriksKienhuesBromme2015METI}.

The expertise-rating interaction was reported as
\(F(1.76,240.01)=100.62,p<.001\). The integrity and benevolence
interactions were \(F(2,272)=341.28,p<.001\) and
\(F(2,272)=341.55,p<.001\), respectively. Planned contrasts
supported greater differences on each targeted scale than on its
comparisons; the integrity-versus-benevolence contrast was smaller than
the integrity-versus-expertise contrast. Corrected degrees of freedom,
within-person dependence, nonnormality concerns and the authors'
effect-size emphasis belong to these results. The main article does
not print numerical confidence intervals for the selected contrasts
\cite[pp.~11--15]{HendriksKienhuesBromme2015METI}.

The studies support a differentiated measurement of source impressions
in these tasks. Study~2 repeats a broad factor pattern in another sample;
Study~3 tests sensitivity to descriptions. They do not measure whether
participants subsequently followed advice, chose information sources,
detected false claims or improved real-world decisions. The authors
explicitly identify the perception-to-actual-trust link as future work.
These studies are not an independent replication across languages and
institutions, and they do not certify a particular expert's reliability
\cite[pp.~16--17]{HendriksKienhuesBromme2015METI}.

\section{Motivation, knowledge and disagreement}
\label{sec:p18-r143-motivations}

The source next asks how goals, identities, prior beliefs and information
selection affect content/source judgments. Selective exposure and selective
avoidance are distinct: seeking agreement need not mean avoiding every
disagreeing source. Its Garrett (2009) report gives stronger preference for
confirming information than avoidance of disconfirming information, at
that literature's scope, not a universal account of all online behavior.

The reviewed gamers example concerns disparagement of findings about
violent games and perceived stigmatization, potentially generalizing to
research on games as a whole. Other reports concern US partisan differences
across sixteen socioscientific issues, free-market beliefs and climate
attribution, and differing acceptance of evolution or vaccination claims.
Those examples do not identify the motives of a particular dissenting
person. Nor do they establish that party membership or an unpopular
conclusion makes every reason offered by its holder invalid. The chapter
provides neither one common participant record nor effect/uncertainty
estimates for this collection \cite[pp.~771--772]{BrommeGierth2021PublicScience}.

Note~7 adds an important objection: reasoning can be motivated in general,
and disagreement can concern the criteria for what counts as evidence,
not only intentional rejection of known evidence. The chapter also
acknowledges disagreement about whether clusters of political, religious
and ideological stances should be called cultures. Motivational explanation
is a hypothesis about a process, not a substitute for assessing the
evidence and arguments actually offered.

The distinction between factual claims and public choices is especially
important. A debate about whether emissions affect temperature differs
from a debate about who bears the costs of a response. Distributional
interests can affect factual beliefs, but scientific evidence alone does
not fix acceptable costs, fairness or legitimate authority. The chapter's
reported wording effect, changing a survey phrase from global warming to
climate change, concerns measured US partisan responses. It is not a
general instruction for manipulating public agreement or a change in
the underlying physical evidence \cite[pp.~771--772 and note~7]{BrommeGierth2021PublicScience}.

The authors criticize a crude enlightenment heuristic: more scientific
information does not automatically yield greater acceptance and support.
They report a small positive knowledge/general-trust association but
more polarized topic-specific beliefs in some educated or knowledgeable
groups. General trust, topic belief, schooling, knowledge and numeracy
are different measures; none is a randomized increase in knowledge merely
because it correlates with another outcome. The chapter supplies no
single synthesis estimate or original participant/measurement record
for the selected polarization comparisons.

This is not a proof that education is useless or counterproductive.
The source confines the concerning findings to contested topics and
specified strong belief patterns, and retains examples of effective
counterargument, including a mechanistic climate explanation. That
example remains a reviewed possibility for attitude change, not a
newly verified universal intervention effect. Assessing an explanation's
correctness and measuring its persuasiveness are separate tasks
\cite[p.~772]{BrommeGierth2021PublicScience}.

\section{Science as an epistemic social system}
\label{sec:p18-r143-stage}

The play/stage distinction expresses the source's conclusion. Public
understanding is one domain in which actual reasoning and normative
standards can differ: it is one play. It is also part of the societal
setting through which people obtain the background knowledge used
in many other domains: it is a stage. This is a philosophical organizing
analogy, not a quantitative causal model of society.

The authors argue that judging scientific institutions can be less bounded
than judging every specialist claim directly. Their relevant institutional
knowledge includes domain consensus, affiliations, vested interests,
reasons for genuine disagreement, expert limitations and failure. Citizens
can relate some criteria to experience with other social institutions.
The source uses underestimated climate-science consensus despite some
mechanistic understanding as an example of why knowledge about science
as an institution differs from knowledge of a technical topic.
Such knowledge itself requires learning; the argument does not assume
that every citizen already possesses it or can identify every deceptive
institution. Motivation and misleading information still threaten these
judgments \cite[pp.~772--773]{BrommeGierth2021PublicScience}.

Consensus is relevant evidence about what the pertinent expert community
currently accepts, not deductive truth or a count of internet repetitions.
Its value depends on competence, independence, evidence and how dissent
was handled. Note~6 distinguishes a citizen's reliance on consensus from
the scientific community's own need to generate and scrutinize reasons.
Scientists cannot discharge their inquiry responsibilities merely by
deferring collectively to the consensus whose warrant they must examine.

Note~8 compares representative democracy's assumption that citizens can
judge whom to trust without personally solving every political problem.
That comparison does not prove all voters equally competent, make expert
accuracy sufficient for democratic legitimacy, or transfer scientific
authority into authority to choose everyone's values. Note~9 identifies
educational approaches emphasizing internet source assessment; neither
note supplies an unrestricted public-decision prescription
\cite[p.~773, notes~6, 8--9]{BrommeGierth2021PublicScience}.

\section{Using testimony without substituting trust for evidence}
\label{sec:p18-r143-use}

\begin{hbmethod}{def:p18-r143-source-assessment}{Assess a scientific claim and its source at the available scope}
Identify the actual question: explaining a phenomenon, accepting a factual
claim, estimating an uncertain outcome, or recommending an action.
Separate what is comprehensible from what can be independently verified.
Record the claim's conditions, uncertainty and actual supporting source.
Check relevant expertise and pertinence, scientific conduct, disclosed
interests, institutional safeguards and the basis of any claimed consensus.
Trace common provenance before treating repeated statements as independent.

Use what technical evidence is genuinely assessable; use source assessment
where specialized justification must be entrusted to others. Examine
documented objections and distinguish a challenged method from contrary
evidence for the conclusion. Retain uncertainty rather than choosing a
claim merely because it is fluent, familiar or ideologically agreeable.
When an action is proposed, separately assess consequences, alternatives,
values, rights and legitimate decision procedures. Revisit reliance when
relevant evidence or source conduct changes.
\end{hbmethod}

This method connects to the actual expert-opinion questions in
\ref{def:p18-expert-opinion-scheme}: expertise, field fit, exact opinion,
trustworthiness, consistency with relevant peers and backup evidence.
It does not replace those questions with a METI score. Witness testimony
under \ref{def:p18-witness-testimony-scheme} additionally needs opportunity
to observe; scientific inference and reporting are not always direct
eyewitness observation. Chapter~\ref{sec:p18-rationality-5-5} separates
being persuaded, evidential strength and dialogue obligations. An acceptable
exchange or a favorable source impression need not yield a true conclusion.

Chapter~\ref{sec:p18-rationality-10-1} distinguishes epistemology in groups
from beliefs attributed to a group, and coherence/truth from legitimacy.
The present source concerns a citizen's dependence on expert knowledge,
not a completeness theorem for opinion aggregation. Repeated reports
from one origin do not supply the independence assumptions of a voting
or pooling model; corporate scientific acceptance is not identical to
every researcher's private belief or to formal common knowledge.

\begin{hbexample}{def:p18-r143-worked-contrast}{A bounded institutional assessment and its failures}
Consider a hypothetical municipality assessing two technical reports
about a proposed carbon-storage project. One report comes from a relevant
research team, states its measurement limits and makes its review process
and affiliations inspectable. A second report presents a confident
opinion by a prominent but unrelated specialist and repeats an undisclosed
source. A citizen unable to reproduce the specialist modeling can still
have reasons to give the first report greater provisional weight.
The comparison concerns relevance, provenance and accountable conduct,
not certainty conferred by prestige or a guarantee that the first model
is correct.

A legitimate objection might identify a failed measurement, an inapplicable
assumption or concealed contrary findings; it must be assessed rather than
dismissed because it challenges an institution. An unsupported accusation
of corruption does not itself establish the opposite technical conclusion.
Even a well-supported engineering assessment leaves questions about costs,
distribution of risk, alternatives and authorization. The trustworthiness
of the technical team does not decide those questions for the municipality.

Two failures are claiming personal mastery because the report is easy
to read, and counting repeated copies of the same assertion as independent
consensus. A third is assuming that honest intentions make an unrelated
specialist competent. These are constructed contrasts implementing the
chapter's distinctions, not observed municipal decisions or empirical
success rates.
\end{hbexample}

\begin{proofboundary}{def:p18-r143-boundary}{Warranted reliance, not unconditional obedience}
The complete 2021 chapter supplies a philosophical and review-level
account of why bounded understanding makes source evaluation important.
Its motivational and knowledge examples retain their original review
standing. The bounded 2015 primary record supports measurement of specific
source impressions, not objective reliability, actual advice following,
or generalized civic competence. No universal public-opinion effect,
deference rule, education failure or political prescription follows.

Rational public dependence on science requires both openness to evidence
and justified attention to the social production of knowledge. Limited
understanding does not make all opinions equally warranted, and expertise
does not settle every public duty. That is the chapter's practical
epistemic distinction, not a new formal system with an unwritten proof.
\end{proofboundary}
